\section{Intro}

Dans un contexte de changement climatique lié aux activités humaines \cite{IPCC2021SPM}, même le pape s'inquiète \cite{leoXIV2026MagnificaHumanitas}

Le rapport « Energy and AI » publié par l’IEA en avril 2025 prévoit, dans son scénario de référence, la demande mondiale en énergie pour les centres de données doubler et passer à environ 945 TWh d’ici 2030 — soit un niveau comparable à la consommation annuelle d’électricité du Japon —, et qui devrait atteindre environ 1 200 TWh d’ici 2035 \cite{iea2025energy}.

Les grandes entreprises, PME, gouvernements, institutions, collectivités, associations et autres organisations doivent contrôler et plafonner leur empreinte environnementale. Des conseils concrets et réels pour la réduction de cet impact sont nécessaires. Mais avant cela, il faut chercher à le comprendre et le mesurer.
Des organismes comme les Nations Unies ont récemment publié un rapport sur l'impact environnemental au niveau des émissions carbones, de la consommation d'eau et de l'utilisation de l'espace \cite{aczel2026environmental}.

De mon côté, en cherchant à réduire cet impact environnemental, j’ai pu développer une méthode améliorant l'efficacité de l'IA générative.
